\section{Introduction}\label{sec:intro}
This lab builds a proportional controller for a simulated kitticopter and tests it on the real
system. The kitticopter moves vertically under rotor thrust, aerodynamic drag and gravity, and
its parameters are unique to each student, so the model cannot be taken from the brief and has
to be measured.

The work runs over two lab sessions. The first identifies the system: the hover voltage is
found, a step test is run at the input, and the resulting response gives the plant gain $A$, the
time constant $T$ and the sensor constant $k$. The second designs a proportional controller from
those parameters, builds it on Veroboard from op-amp stages, and demonstrates it tracking a
position setpoint.

The controller has to track position setpoints to within $8$ percent, keep overshoot below $30$
percent, and hold both of those in the face of $10$ percent uncertainty in the aerodynamic
constant and $10$ percent tolerance in the components used to build it. This report derives the
model, reports the measurements and the parameters extracted from them, shows how the gain was
chosen against those four specifications, and evaluates what the built controller actually
achieved.
